\subsection{Side-specific interval proximal update}\label{sec:asymmetric}
The side-specific projection path evaluates the separable penalty
\eqref{eq:side-price} inside an original-row OSQP-style iteration. Its narrow
implementation property is that lower and upper linear/quadratic coefficients
enter the row operation without adding independent departure variables to the
principal system. Unequal side costs already have established formulations,
including the DAQP interface. The question tested by the original experiment
is computational cost against equivalent lifting, not whether unequal prices
can be modeled for the first time.

For a temporary scalar value $v$ and positive numerical metric $\rho$, the row
operation solves
\begin{equation}\label{eq:asym-subproblem}
 z^+=\argmin_z\{g_i(z)+\tfrac\rho2(z-v)^2\}.
\end{equation}
It has the closed form
\begin{equation}\label{eq:asym-prox}
z^+=\begin{cases}
\ell_i-\dfrac{[\rho(\ell_i-v)-a_{L,i}]_+}{\rho+b_{L,i}},&v<\ell_i,\\[5pt]
v,&\ell_i\leq v\leq u_i,\\[3pt]
u_i+\dfrac{[\rho(v-u_i)-a_{U,i}]_+}{\rho+b_{U,i}},&v>u_i.
\end{cases}
\end{equation}
An absent side is omitted. A hard row uses exact interval projection. The
positive denominator follows from $\rho>0$ and nonnegative quadratic price;
the formula remains defined for a pure linear or a zero-price soft side.

To derive the upper branch, write $d=z-u_i>0$. Differentiating the exterior
quadratic yields $(\rho+b_{U,i})d=\rho(v-u_i)-a_{U,i}$. If its right-hand side
is nonpositive, the exterior optimum is inadmissible and the endpoint satisfies
the subgradient condition. Otherwise division gives the positive departure.
The lower branch follows by reflecting the displacement and using the lower
coefficients. Inside the interval the price is zero and $z=v$ minimizes the
proximal displacement. Strong convexity of \eqref{eq:asym-subproblem} establishes
uniqueness, including at the interval endpoints.

This scalar derivation is an application of established proximal convex
analysis~\cite{prox,moreau,rockafellarppa}. It is not a proof of arbitrary adaptive
ADMM behavior, a finite-precision error theorem or a global priority claim.
Within a fixed positive metric and an attained primal and dual solution, the
underlying splitting framework supplies the relevant convergence argument.
Unit conversion must transform the side prices, row coefficients and dual
coordinates consistently. A positive absolute tolerance in physical departure
units cannot by itself establish covariance across metres and centimetres.

Algorithm~\ref{alg:asym} identifies the actual native procedure. The separate
engine owns the C/C++ update, Python bindings, scaling and objective recovery.
The product consumes the public SDK in an isolated profile. Grouped and Huber
costs do not become supported merely because their scalar notation resembles
an elastic price. The separable adapter refuses them; the mixed C13 model uses
the published generalized projection or exact lifting instead.

\begin{algorithm}[htbp]
\caption{Apply side-specific scalar proximal updates in the original interval rows while preserving hard policy and independently recovering the requested objective and dual conditions.}\label{alg:asym}
\begin{algorithmic}[1]
\Require Admitted separable original problem, side prices and positive row metric
\State Prepare the fixed-shape original-row native workspace
\For{each row-update operation}
\State Project hard rows; apply \eqref{eq:asym-prox} to each admitted soft row
\State Continue the native primal/dual splitting update
\EndFor
\State Recover original $x,y$ and physical lower/upper departures
\State Evaluate the original objective and method-specific termination criteria
\State Compare accepted outputs to the independent equivalent lift
\end{algorithmic}
\end{algorithm}

The implementation must preserve more than the scalar minimizer. Original
objective evaluation, the dual conjugate, scaling, recession behavior and
fixed-shape repeated updates have independent controls. A stationary vector
can still violate the soft subgradient relation. The diagnostic proximal test
$z=\prox_g^{H_0}(z+H_0^{-1}y)$ checks that relation at the original row value rather than
selecting a derivative arbitrarily at a kink. Here $H_0$ is an explicitly
positive diagonal diagnostic metric; its numeric identity applies only in
canonical diagnostic coordinates. Signed slack cancellation is
prevented by separate lower and upper departure arrays.

The original held-out endpoint does not establish the registered computational
gain. Its median relative reduction is \OriginalEstimate\%, with a confidence
interval crossing zero. Correctness and retained original dimensions therefore
coexist with a negative timing conclusion. A favorable selected family, one
small scalar matrix or a quick warm update cannot replace that primary result.
The implementation remains a complete executable method within its supported
domain; its scientific limitations are stated through the full comparison,
not hidden by relabeling the failed prediction as a different acceptance rule.
Endpoint behavior supplies useful concrete controls. A sufficiently large
linear price relative to one temporary proximal displacement maps that
displacement exactly to the preferred endpoint. This is a statement about
one proximal subproblem, not a guarantee that a finite price enforces the
original row for every possible model or iteration. Setting the quadratic
price to zero leaves a genuine one-sided soft-threshold operation. Setting
both prices to zero leaves a free preference, while an explicit hard mask
still projects exactly. Confusing these three cases changes the optimization
policy even though each resulting program can remain convex and solvable.
The scalar analytic cases and independently recovered original dual signs
test these distinctions without relying on matching backend status integers
or comparing an auxiliary representation to itself.
